{\large\textbf{Cat eyes}}

You may have noticed that in darkness, when a cat is within the light beam of a
headlamp, its eyes appear very bright, see the photo below (left). This
phenomenon can be modelled by a lens setup, see the photo on right, and the
diagram beneath the photos.

\begin{center}\includegraphics[width=0.89\linewidth]{figures/EuPhO_2020_Q3_fig1.jpg}\end{center}

\begin{center}\includegraphics[width=0.89\linewidth]{figures/EuPhO_2020_Q3_fig2.png}\end{center}

The photo on right was taken by a digital single-lens reflex camera. The light
intensity at the camera sensor pixels marked by a red line (in the photo) is
shown in the graph below: the log base 10 of the light intensity (measured as
the number of photons caught by each pixel) is plotted against the
$x$-coordinate, with the pixels' side length serving as the unit length.

\begin{center}\includegraphics[width=0.89\linewidth]{figures/EuPhO_2020_Q3_fig3.png}\end{center}

The lens modelling cat eyes can be treated as an ideal thin lens of focal
length $f = 55\ mm$ and diameter $D = 39\ mm$; however, you should keep in mind
that the given graph shows real measurement data, and the lens has certain
non-ideal features. Most importantly, partial reflections of brightly lit areas
from the lens surfaces may decrease the contrast: dark areas seen through the
lens appear less dark than they actually are; this effect can be neglected for
the camera lens, but not so for the lens serving as a model of a cat's eye.

Based on the given data, estimate (with the accuracy of \textit{ca} $20\%$ )
the distance $h$ between the axis of the camera and the axis of the lamp (which
can be considered as a point source) if the distance of the camera from the
paper sheet was $L = 4.8\ m$.
